%=======================================================================================
% METHODOLOGY CHAPTER -- 2026-08-11, 22:21
% Reconstructed from the implemented pipeline, not from earlier plans. Every claim below is
% checked against the actual code (file paths given inline) or explicitly marked
% "% TODO: verify" / "% TODO: citation" where it could not be confirmed this pass. No
% performance results, interpreted associations, or model-comparison verdicts appear here --
% see the Results chapter (and TEMP_results/README.md, an interim index of completed runs) for
% those. Detailed hyperparameter grids, full VIF/correlation outputs, and per-variable
% diagnostics are pointed at the Appendix rather than repeated here -- see
% documentation/methodlogy_env_setpick.md and documentation/env_feature_sets_manifest.csv for
% the source material an Appendix chapter should draw on.
%=======================================================================================

\chapter{Methodology}
\label{ch:methodology}

This chapter describes how the raw plot data were converted into the results used to answer the
project's three research questions: which model best predicts top height (RQ1), what is
associated with departure from a single shared Chapman-Richards curve (RQ2), and what is
associated with plot-specific deviation in estimated growth trajectory (RQ3). The chapter is
organised by method, not by the order experiments were run: common design decisions first
(\S\ref{sec:common-design}), then environmental feature-set construction
(\S\ref{sec:feature-sets}), then the three RQs in turn (\S\ref{sec:rq1}--\S\ref{sec:rq3}), and
finally reproducibility and appendix pointers (\S\ref{sec:reproducibility}).

\section{Common experimental and evaluation design}
\label{sec:common-design}

\subsection{Cohorts}

Two overlapping cohorts are used throughout: \textbf{4survey} (four repeated LiDAR
surveys per plot) and \textbf{6survey} (six repeated surveys, a strict subset of 4survey's plots
with a longer, denser survey history but a smaller training set). 4survey is the primary cohort
for every research question; 6survey is run in parallel as a secondary check wherever the sample
size permits it. Dataset-level detail (plot counts, survey years, coverage) is given in the Data
chapter and is not repeated here.

\subsection{Splits and cross-validation}
\label{sec:splits}

Four split types are used, defined in \texttt{models/common/splits.py}, each answering a
different generalisation question:

\begin{itemize}
\setlength\itemsep{0pt}
\item \textbf{Spatial block} (primary). Whole forestry compartments (\texttt{cpmt}, the block
  identifier column) are assigned to train/validation/test so that no compartment appears on both
  sides of a split -- a plot-level or random split would let a model see a near-identical
  neighbouring plot during training, overstating how well it generalises to genuinely new
  ground. Compartments are chosen by row-count-balanced random assignment (60/20/20
  train/validation/test, seed 42), with a 60\,m buffer around every validation/test plot: any
  \emph{training} plot within 60\,m of a held-out plot is excluded from training (validation/test
  plots are never removed). The 60\,m figure is set from typical plot spacing, not from a
  measured spatial-autocorrelation range -- the shared-curve residual's own semivariogram range is
  approximately 3{,}956\,m, far larger than the buffer. Leakage protection at that larger scale
  comes from holding out whole compartments, not from the buffer alone
  \citep{roberts2017}. % TODO: citation -- Roberts et al. 2017, cited in splits.py's own comments for the spatial-block-CV rationale; no .bib entry exists yet.
\item \textbf{Spatial block, pooled $k$-fold} (stronger confirmation of the single spatial-block
  split). Compartments are assigned to 5 row-count-balanced folds; each fold in turn is held out
  as test, the next as validation, the remainder as train, so every compartment is evaluated
  exactly once. Predictions from all 5 folds are pooled into one population-wide evaluation set
  rather than averaged as 5 separate summary statistics.
\item \textbf{Temporal} (secondary, harder test). Whole survey \emph{years} are assigned to
  train/validation/test instead of compartments: 4survey trains on 2008 and 2012, validates on
  2021, tests on 2023; 6survey trains on 2002/2006/2008/2012, validates on 2021, tests on 2023.
  This asks whether a model trained on earlier growth still predicts a real future survey it
  never saw, a different and harder question than spatial generalisation. A narrow-gap variant
  (training years closer to the test year) is also defined for a secondary check; its validation
  year is deliberately the earliest available year, not a middle year, because holding out a
  middle year left the PINN's trajectory-pair loss with no pairable plots in one smoke test.
\item \textbf{Plot level} (sanity floor only, not a generalisation claim). Individual plots are
  split at random (60/20/20, seed 42) with no spatial holdout at all -- kept only as an easy-case
  reference, since a held-out plot's nearest neighbour is typically a training plot metres away.
\end{itemize}

Every model in this project reads the same split assignment for a given (cohort, split type,
seed) combination, so results are always compared on identical train/validation/test membership.
\texttt{cpmt}/\texttt{blk}/\texttt{scpt} (the compartment/block/sub-compartment identifier
columns) are asserted to never appear in any model's own feature list, enforced automatically at
import time.

\subsection{Feature scaling, missing values, and encoding}

Continuous features are standardised (zero mean, unit variance) using a scaler fitted on the
training partition only, then applied unchanged to validation/test -- avoiding any information
about the held-out rows leaking into the scaling itself. Age is scaled with its own separate
scaler, kept apart from other continuous features so its training-split standard deviation can be
read back out directly for the physics loss's unit correction (\S\ref{sec:rq1-pinn}). Unordered
categorical variables (three CEH soil-class rasters) are one-hot encoded, never treated as an
ordered numeric scale. Rows with a missing value in any requested feature are dropped whole-plot
(every survey year of that plot, not just the affected row), so no model ever sees a partial
trajectory for a plot it has otherwise seen.

\subsection{Evaluation metrics}
\label{sec:metrics}

Every model is scored on the same metric set, computed from residuals $e_i = H_i - \hat H_i$
(observed minus predicted height) over $N$ rows:
\begin{align*}
\text{MAE} &= \frac{1}{N}\sum_i |e_i|, &
\text{RMSE} &= \sqrt{\frac{1}{N}\sum_i e_i^2}, \\
R^2 &= 1 - \frac{\sum_i e_i^2}{\sum_i (H_i - \bar H)^2}, &
\text{Bias} &= \frac{1}{N}\sum_i e_i,
\end{align*}
where $\bar H$ is the mean observed height in that evaluation set; a positive bias indicates
systematic under-prediction. Metrics are also broken down by five age bands ($<30$, 30--40,
40--60, 60--80, $80+$ years) to check whether an aggregate score conceals age-specific weaknesses;
the maturity gate (plot-level \texttt{Age}~$\geq$~30 at the 2023 reference survey) is applied at
plot level, so individual rows below 30 can still appear in a trajectory, and a poor score in the
$<30$ band reflects a small, non-representative sample rather than necessarily worse prediction.

\subsection{Spatial residual diagnostics}

Moran's I is computed on model residuals (not on raw features, and not as a variable-selection
criterion) to test whether spatial structure remains unexplained after fitting. This is a
diagnostic applied after a model is fitted, answering a different question from environmental
feature selection (\S\ref{sec:feature-sets}); it is not used anywhere to decide which variables
enter a model.

\section{Environmental feature-set construction}
\label{sec:feature-sets}

Three questions each need their own screened list of candidate environmental/management
variables: RQ1's predictive tiers, RQ2's shared-curve attribution set, and RQ3's plot-level
attribution set. All three follow the same five-stage recipe, differing only in candidate pool
and target column (each explicitly parameterised, so a ranking can never be accidentally computed
against the wrong target). Full per-variable diagnostics, VIF tables, and the complete final
column lists are given in Appendix~\ref{app:feature-sets}
(source: \texttt{documentation/methodlogy\_env\_setpick.md},
\texttt{documentation/env\_feature\_sets\_manifest.csv}); this section states the rules only.

\begin{enumerate}
\setlength\itemsep{0pt}
\item \textbf{Baseline.} Four stand/management columns (canopy cover, time since thinning, a
  paired missingness flag, and recent-thinning fraction) are present unconditionally in every
  set. A fifth candidate, thinning fraction, was confirmed to be an exact linear complement of
  the missingness flag and is not included separately.
\item \textbf{Deduplication.} Deterministic duplicates (a variable that is a documented closed-form
  function of another candidate already in the pool) are dropped outright; near-exact empirical
  duplicates (Spearman $|\rho|\geq0.95$) are reduced to one representative per pair, preferring an
  externally measured variable over a derived one.
\item \textbf{Importance ranking.} Three signals are computed against each avenue's own target and
  combined by rank-averaging, not by any one signal alone: univariate Spearman $|\rho|$; XGBoost
  permutation importance; and XGBoost drop-column ablation (the held-out $R^2$ change from
  removing one variable, given every other candidate already in the model). A single ranking
  scheme has a known blind spot -- correlation ignores context, permutation is biased under
  correlated predictors, ablation depends on one model/split -- so a variable must look important
  under more than one lens to rank highly overall.
\item \textbf{Variance-inflation screening.} Iterative VIF reduction (threshold 5.0
  \citep{dormann2013}) is applied on top of ranking/gating, with the four baseline columns
  protected from removal (they remain in the design matrix, affecting every other column's VIF,
  but are never themselves dropped).
\item \textbf{Final sets.} Set1 = baseline only. Set2 = baseline + the top 10 candidates by
  combined rank, itself VIF-screened with backfill from the next-ranked candidate (so a fixed-size
  top-10 pick cannot silently include a redundant pair). Set3 = baseline + terrain/wind
  candidates in the upper half of their own category by combined rank, VIF-screened. Set4 = Set3
  plus every other category's upper-half candidates, VIF-screened. Set5 = baseline plus every
  deduplicated candidate, VIF-screened, no rank filter.
\end{enumerate}

\begin{table}[!htbp]
\centering
\scriptsize
\renewcommand{\arraystretch}{1.2}
\begin{tabular}{p{2.3cm}p{3.6cm}p{1.6cm}p{2cm}p{4cm}}
\toprule
Feature set family & Feeds & Model comparison & Attribution & Purpose \\
\midrule
RQ1 Set1--Set5 & \texttt{dnn\_env\_terrain}, \texttt{pinn\_env\_terrain*} & Primary & No & DNN-vs-PINN
comparison under progressively richer terrain/wind/climate input; no coefficient-based model
reads these lists. \\
RQ2 Set1--Set5 & Elastic Net, XGBoost, NLME & Secondary & Primary & Which conditions align with
persistent departure from the shared curve. \\
RQ3 Set1--Set5 & Elastic Net, XGBoost, GNNWR & Secondary & Primary & Which conditions align with
plot-specific trajectory deviation; GNNWR additionally tests whether that association is
spatially varying. \\
\bottomrule
\end{tabular}
\caption{Purpose of the final feature sets by research question. Full column lists: Appendix~\ref{app:feature-sets}.}
\label{tab:feature-set-purpose}
\end{table}

% FIGURE PLACEHOLDER (pipeline overview)
% Caption: "From plots to results: cohorts and candidate environmental features feed a shared
% five-stage screening process (dedup, importance ranking, VIF), producing Set1-5 per research
% question. RQ1's sets feed the predictive DNN/PINN branch under spatial/temporal/pooled-k-fold
% splits; RQ2/RQ3's sets feed the two attribution branches (shared-curve residual; plot-specific
% trajectory), each producing Elastic Net/XGBoost/(NLME or GNNWR) fits evaluated under spatial
% block cross-validation."
% Content: a single left-to-right flow diagram. Boxes: [Raw plot/survey data + candidate
% environmental variables] -> [Cohort selection: 4survey / 6survey] -> [Feature screening:
% dedup -> combined-rank importance -> VIF, per RQ] -> [Set1-5, per RQ] -> three parallel
% branches: (a) "RQ1: predictive models" -> CR/average-by-age/Linear/RF/XGBoost baselines +
% DNN/PINN variants -> spatial/temporal/pooled-kfold evaluation; (b) "RQ2: shared-curve
% attribution" -> mean_cr_residual -> Elastic Net/XGBoost/NLME; (c) "RQ3: plot-specific
% attribution" -> local_y_max_difference -> Elastic Net/XGBoost/GNNWR. All three branches
% converge on a final "Results" box.
\begin{figure}[!htbp]
\centering
\fbox{\parbox{0.9\textwidth}{\centering\vspace{4cm}FIGURE PLACEHOLDER --- pipeline overview, see
comment above source for exact content\vspace{4cm}}}
\caption{Overview of the full pipeline from raw data to the three research questions' results.}
\label{fig:pipeline-overview}
\end{figure}

\section{Top-height prediction models}
\label{sec:rq1}
\label{sec:rq1-pinn}

Target: \texttt{elev\_percentile\_95th} (95th-percentile LiDAR-derived top height, per plot per
survey). Six model families are compared, sharing splits, metrics, and (where applicable) network
architecture, so any performance difference is attributable to the model itself rather than to a
confounded training difference.

\subsection{Baselines}

\textbf{Chapman-Richards (CR) curve.} A three-parameter curve
\begin{equation}
H(a) = y_{\max}\left(1-e^{-ka}\right)^{p}
\label{eq:cr}
\end{equation}
is fitted by nonlinear least squares (\texttt{scipy.optimize.curve\_fit}, five restarts, lowest
sum-of-squared-error kept), where $H(a)$ is predicted top height at age $a$ (years), $y_{\max}$ is
the asymptotic maximum height, $k$ controls growth rate, and $p$ is a shape parameter. This is a
single \emph{pooled}, population-average curve -- one fit per cohort, not per plot or per
compartment -- fitted on the training partition of the split in use (never on validation/test
rows), refitted separately per fold under the pooled spatial-block $k$-fold split so that a given
fold's held-out compartments never influence its own frozen curve. The pooled curve's fitted
$(y_{\max},k,p)$ are treated as constants downstream (\S\ref{sec:rq1-pinn-loss}), never as
trainable parameters.

\textbf{Average-by-age.} A lookup table of mean observed height per integer age year, computed on
the same training rows as the CR fit; an unseen age falls back to the overall training mean
height.

\textbf{Linear, Random Forest, and XGBoost baselines.} Conventional supervised baselines fit
directly on the no-environment feature set, evaluated under the same splits as every other model.
XGBoost's hyperparameters (\texttt{max\_depth}~$\in\{2,3,4,6\}$,
\texttt{reg\_lambda}~$\in\{1,5,20\}$, \texttt{n\_estimators}~=~500 fixed with early stopping,
\texttt{learning\_rate}~=~0.1 fixed) are selected by validation-set $R^2$ over this 12-combination
grid; full grid and selection results: Appendix~\ref{app:hyperparameters}. %
TODO: verify -- confirm whether the environmental-feature-set XGBoost fits reported in RQ1's
sweep used this same grid/selection procedure or a separate script's own configuration.

\subsection{DNN}

A single shared architecture (\texttt{NoEnvNetwork}) is used for every DNN and PINN variant in
this project: three hidden layers of 128 units each, leaky-ReLU activation, age concatenated with
every other input feature. Reusing one architecture class throughout means that any difference
between the DNN and PINN's results is attributable to the loss function, not to a different
network. The DNN's loss is plain mean squared error,
\begin{equation}
\mathcal{L}_{\text{data}} = \frac{1}{N}\sum_{i=1}^{N}\left(\hat H_i - H_i\right)^2,
\label{eq:data-loss}
\end{equation}
where $\hat H_i$ is the network's predicted height for row $i$ and $H_i$ is the observed height,
plus an L1 weight penalty ($\lambda_{L1}=10^{-5}$). \texttt{dnn\_env\_terrain} is architecturally
identical to \texttt{dnn\_noenv}; its terrain/wind block is simply concatenated into the same
"other features" input, widening that vector rather than adding a second network.

\subsection{PINN variants}

Three physics-informed variants share the DNN's exact network class, differing only in loss
terms and, for the environment-conditioned variants, in an added small sub-network:

\begin{itemize}
\setlength\itemsep{0pt}
\item \texttt{pinn\_noenv} -- no environmental input; the CR curve's $(y_{\max},k,p)$ are the
  single pooled, frozen cohort constants from Equation~\ref{eq:cr}.
\item \texttt{pinn\_env\_terrain} -- terrain/wind features feed a small sub-network
  (\texttt{YMaxSubNetwork}, 16 hidden units) that outputs a per-plot \emph{additive} adjustment to
  the pooled $y_{\max}$:
  \begin{equation}
  y_{\max}^{(i)} = y_{\max} + f_{y}\!\left(\mathbf{x}^{\text{terrain}}_i\right),
  \label{eq:ymax-adjust}
  \end{equation}
  where $f_y$ is the sub-network and $\mathbf{x}^{\text{terrain}}_i$ is plot $i$'s terrain/wind
  feature vector. $k$ and $p$ remain the pooled global constants. This sub-network never sees age
  or the no-environment features, and its output is initialised near zero so training starts from
  the known-good pooled curve rather than an arbitrary value. Terrain/wind conditions only the
  growth \emph{ceiling} in this design, not the trajectory shape -- matching the
  asymptote-varies, shape-fixed framing of site-index modelling
  \citep{socha2021}. % TODO: citation -- Socha et al. 2021 (ADA/GADA), no .bib entry yet.
\item \texttt{pinn\_env\_terrain\_k} -- adds a second sub-network producing a per-plot
  \emph{multiplicative}, log-space adjustment to $k$:
  \begin{equation}
  k^{(i)} = k \cdot \exp\!\left(f_{k}\!\left(\mathbf{x}^{\text{terrain}}_i\right)\right),
  \label{eq:k-adjust}
  \end{equation}
  again initialised near zero so $k^{(i)}\approx k$ at the start of training. $p$ stays a pooled
  global constant in every variant; no variant makes the curve's shape parameter site-specific.
\end{itemize}

\paragraph{Loss function.} Every PINN variant optimises the same three-term objective, using
$H'(a)$, the CR curve's analytical derivative,
\begin{equation}
H'(a) = y_{\max}\,p\,k\,e^{-ka}\left(1-e^{-ka}\right)^{p-1}.
\label{eq:cr-derivative}
\end{equation}
The \emph{physics} term matches the network's own instantaneous growth rate, obtained by automatic
differentiation of its prediction with respect to age, against the analytical rate at that same
age:
\begin{equation}
\mathcal{L}_{\text{phys}} = \frac{1}{N}\sum_{i=1}^{N}\left(\frac{\partial \hat H_i}{\partial a_i} -
H'(a_i)\right)^2.
\label{eq:physics-loss}
\end{equation}
The \emph{trajectory} term is a separate, finite-difference check: for $M$ pairs of real,
consecutive same-plot survey observations (both endpoints drawn only from that plot's own
training-labelled rows), the network's own predicted height \emph{change} between the pair must
match the analytical rate at the pair's mid-age,
\begin{equation}
\mathcal{L}_{\text{traj}} = \frac{1}{M}\sum_{j=1}^{M}\left(\frac{\hat H_{j,\text{later}} -
\hat H_{j,\text{earlier}}}{\Delta a_j} - H'(\bar a_j)\right)^2,
\label{eq:trajectory-loss}
\end{equation}
where $\Delta a_j$ is the elapsed time between the pair's two surveys and $\bar a_j$ is their
mid-age. The two terms are deliberately distinct checks: an instantaneous derivative at one point
versus a measured change across two real observations. The full training objective is
\begin{equation}
\mathcal{L} = \mathcal{L}_{\text{data}} + \lambda_{\text{phys}}\mathcal{L}_{\text{phys}} +
\lambda_{\text{traj}}\mathcal{L}_{\text{traj}} + \lambda_{L1}\sum_{\theta}|\theta|,
\label{eq:total-loss}
\end{equation}
with $\lambda_{\text{phys}}=\lambda_{\text{traj}}=1.0$ as the default weighting (a physics-weight
ablation over $\{0,1,2\}$ was also run; settings and full results: Results chapter, not repeated
here). $(y_{\max},k,p)$ are always plain constants at every point in this objective, never
tensors that accumulate a gradient -- training updates only the network (and, where present, the
$y_{\max}$/$k$ sub-networks), never the process model itself. A trajectory pair is only formed
from two rows both labelled training under the split in use, so no validation or test survey year
can leak into this loss term.

% FIGURE PLACEHOLDER (model architecture)
% Caption: "Shared DNN/PINN architecture. The plain data loss (DNN, and the PINN's own data
% term) uses the main trunk's output directly. The PINN adds two consistency terms against the
% Chapman-Richards curve's analytical derivative: an instantaneous check via automatic
% differentiation (physics loss) and a finite-difference check across real survey pairs
% (trajectory loss). In the environment-conditioned variants, terrain/wind features feed only
% the small y_max (and, in the _k variant, k) sub-networks -- never the main trunk -- producing
% a per-plot adjustment to the CR curve's asymptote/rate that the physics and trajectory losses
% are then computed against."
% Content: a block diagram. Left: [Age, no-env features] -> main trunk (3x128 NoEnvNetwork) ->
% [predicted height]. Below/beside: [Terrain/wind features] -> small y_max sub-network (16
% units) -> [y_max adjustment, additive] and, dashed/optional for the _k variant, -> k
% sub-network -> [k adjustment, multiplicative, log-space]. A separate branch shows [predicted
% height] and [age] feeding into an autograd box labelled dH/da, compared against a box labelled
% "CR analytical derivative (frozen y_max, k, p)" to produce the physics loss; a second small
% diagram shows two real survey observations of one plot feeding two forward passes, differenced
% and compared to the same CR derivative at mid-age, to produce the trajectory loss.
\begin{figure}[!htbp]
\centering
\fbox{\parbox{0.9\textwidth}{\centering\vspace{4cm}FIGURE PLACEHOLDER --- model architecture, see
comment above source for exact content\vspace{4cm}}}
\caption{Shared DNN/PINN architecture and the two physics-consistency loss terms.}
\label{fig:model-architecture}
\end{figure}

\begin{table}[!htbp]
\centering
\scriptsize
\renewcommand{\arraystretch}{1.2}
\begin{tabular}{p{2.3cm}p{3cm}p{2.6cm}p{4.5cm}}
\toprule
Model & Environmental input & Loss terms & Notes \\
\midrule
CR (baseline) & None & N/A (curve fit) & Pooled per cohort, refit per fold for pooled $k$-fold. \\
Average-by-age & None & N/A (lookup) & Same training rows as CR. \\
Linear / RF / XGBoost & No-environment feature set & Model-specific & XGBoost hyperparameters
selected by validation $R^2$ (Appendix~\ref{app:hyperparameters}). \\
\texttt{dnn\_noenv} & None & Data (Eq.~\ref{eq:data-loss}) + L1 & Shared architecture with every
PINN variant. \\
\texttt{dnn\_env\_terrain} & Terrain/wind, concatenated into main input & Data + L1 & Same network
class as \texttt{dnn\_noenv}, wider input only. \\
\texttt{pinn\_noenv} & None & Data + physics + trajectory + L1 & $y_{\max},k,p$ pooled, frozen. \\
\texttt{pinn\_env\_terrain} & Terrain/wind, $y_{\max}$ sub-network only & Data + physics +
trajectory + L1 & Eq.~\ref{eq:ymax-adjust}; $k,p$ pooled. \\
\texttt{pinn\_env\_terrain\_k} & Terrain/wind, $y_{\max}$ and $k$ sub-networks & Data + physics +
trajectory + L1 & Eqs.~\ref{eq:ymax-adjust}--\ref{eq:k-adjust}; $p$ pooled. \\
\bottomrule
\end{tabular}
\caption{Model family comparison for top-height prediction (RQ1). Full hyperparameters:
Appendix~\ref{app:hyperparameters}.}
\label{tab:rq1-models}
\end{table}

\section{Shared Chapman-Richards residual attribution (RQ2)}
\label{sec:rq2}

This avenue asks which environmental/management conditions align with a plot's persistent
departure from the single shared curve fitted in Equation~\ref{eq:cr}, using only the 4survey
cohort (6survey's 47 compartments were judged too few for a reliable spatial-block attribution
split). For each row $i$ of a plot's own survey history, a residual is computed against the
pooled curve,
\begin{equation}
r_i = H_i - H(a_i),
\label{eq:cr-residual}
\end{equation}
and the target used for attribution is the plot's own mean residual across its $T$ observed
survey rows,
\begin{equation}
\overline{r}_{\text{plot}} = \frac{1}{T}\sum_{t=1}^{T} r_{i,t}.
\label{eq:mean-cr-residual}
\end{equation}
Age is excluded from every RQ2 candidate feature: it is the CR curve's only input and therefore
the sole construction variable of $\overline{r}_{\text{plot}}$, making it circular for this
target. % TODO: verify -- confirm whether any plot contributing to Equation~\ref{eq:mean-cr-residual}'s
average (computed over all of that plot's survey rows) falls in the test partition of the
downstream attribution split below; the target is constructed before that split is drawn, which
is expected rather than a leakage bug, but is not yet explicitly documented as checked.

\subsection{Fitted models}

Three models are fitted on $\overline{r}_{\text{plot}}$ using RQ2's Set1--Set5
(\S\ref{sec:feature-sets}), under the same spatial-block design as \S\ref{sec:splits}:

\begin{itemize}
\setlength\itemsep{0pt}
\item \textbf{Elastic Net} -- continuous features standardised on the training split; unordered
  soil-class variables one-hot encoded; regularisation strength and $\ell_1$/$\ell_2$ mixing
  ratio selected by 5-fold cross-validation on the training partition only (never touching
  validation or test).
\item \textbf{XGBoost} -- same hyperparameter-selection approach as \S\ref{sec:rq1}.
\item \textbf{NLME (two-stage approximation).} A true nonlinear mixed-effects refit of the CR
  curve itself was not implemented (no equivalent Python tooling was judged mature enough for
  this project's timeline). Instead, a two-stage approximation is used: Stage 1 is the already
  frozen, pooled CR curve above; Stage 2 is a \emph{linear} mixed model fitted on the residual
  $\overline{r}_{\text{plot}}$, with the candidate environmental variables as fixed effects and a
  random intercept grouped by forestry compartment,
  \begin{equation}
  \overline{r}_{\text{plot}} = \boldsymbol{\beta}^{\top}\mathbf{x}_{\text{plot}} +
  u_{\text{cpmt}} + \varepsilon_{\text{plot}},
  \label{eq:nlme}
  \end{equation}
  where $\boldsymbol{\beta}$ are fixed-effect coefficients on the standardised feature vector
  $\mathbf{x}_{\text{plot}}$, $u_{\text{cpmt}}\sim\mathcal{N}(0,\sigma_u^2)$ is a compartment-level
  random intercept, and $\varepsilon_{\text{plot}}$ is residual error. Fitted by restricted
  maximum likelihood (REML) for coefficient reporting; a separate variance-explained comparison
  between nested fixed-effects specifications uses maximum likelihood (ML) instead, since REML
  variance components are not directly comparable across models that differ in fixed-effects
  structure \citep{pinheiro2000}. % TODO: citation -- Pinheiro & Bates, no .bib entry yet.
  NLME is fitted on the training partition only and has no held-out test evaluation step; its
  role is coefficient/variance reporting on the fitted training sample, not held-out prediction.
  % TODO: verify -- confirm statsmodels' default random-effects structure (random intercept only,
  no random slope) against the installed version, since no explicit re\_formula is passed in code.
\end{itemize}

\section{Plot-specific trajectory attribution (RQ3)}
\label{sec:rq3}

This avenue asks the same broad question at plot level rather than pooled level: does a plot's
\emph{own} fitted growth ceiling deviate from what a static, external yield-class table would
predict for it, and does that deviation associate with environmental/management conditions.

\subsection{Target construction}

For each plot, a per-plot asymptote $\hat y_{\max}^{\text{plot}}$ is fitted using \emph{only that
plot's own} repeated height/age observations -- never pooled across plots, unlike
Equation~\ref{eq:cr}. Because the curve's shape parameters are held fixed per plot (taken from
its own Forestry Commission yield-class lookup, not refitted), height is linear in the one
remaining unknown, $y_{\max}$, so the fit is closed-form weighted least-squares through the
origin rather than an iterative optimisation. The target is the difference between this
plot-specific fitted asymptote and the asymptote implied by the plot's own static yield-class
table entry,
\begin{equation}
\Delta y_{\max}^{\text{plot}} = \hat y_{\max}^{\text{plot}} - y_{\max}^{\text{yldc}}.
\label{eq:local-ymax-diff}
\end{equation}
As with RQ2, this target is built entirely from height/age survey data and a static external
lookup table, with no environmental variable involved in its construction.

\subsection{Fitted models}

Elastic Net and XGBoost are fitted as in \S\ref{sec:rq2}, on RQ3's own Set1--Set5. A third model,
\textbf{GNNWR} (geographically and neurally weighted regression), additionally allows coefficients
to vary continuously over space rather than assuming one global linear relationship,
\begin{equation}
\hat y_i = \beta_0(s_i) + \sum_j \beta_j(s_i)\,x_{ij},
\label{eq:gnnwr}
\end{equation}
where $s_i$ is plot $i$'s spatial location and each $\beta_j(s_i)$ is a locally varying
coefficient rather than a single global one \citep{du2020}. % TODO: citation -- Du et al. 2020, no .bib entry yet.
This tests whether the attribution result itself is spatially non-stationary, a question the
global Elastic Net/XGBoost fits cannot answer. Gain-based XGBoost importance and SHAP values are
computed only after model fitting, on the full population with no held-out split, purely to
interpret an already-chosen model -- neither is used at any point to decide which variables enter
a feature set (\S\ref{sec:feature-sets}).

% FIGURE PLACEHOLDER (attribution comparison)
% Caption: "The two attribution avenues share the same pooled Chapman-Richards curve as their
% starting point but diverge in what they attribute. RQ2 pools every plot's residual against one
% shared curve and asks what conditions align with a plot's persistent position above/below it.
% RQ3 fits each plot its own asymptote from only that plot's own repeated observations and
% compares it to an external, static yield-class expectation, asking what conditions align with
% plot-specific deviation. Association only in both cases, never causal language."
% Content: two parallel small diagrams side by side under one shared "pooled CR curve" box at
% the top. Left branch (RQ2): [pooled CR curve] -> [per-row residual, Eq. cr-residual] ->
% [mean per plot, Eq. mean-cr-residual] -> [Elastic Net / XGBoost / NLME]. Right branch (RQ3):
% [plot's own repeated observations] -> [closed-form per-plot y_max fit] -> [difference vs.
% static yield-class y_max, Eq. local-ymax-diff] -> [Elastic Net / XGBoost / GNNWR]. A small
% annotation under both branches: "association, not causation."
\begin{figure}[!htbp]
\centering
\fbox{\parbox{0.9\textwidth}{\centering\vspace{4cm}FIGURE PLACEHOLDER --- attribution comparison,
see comment above source for exact content\vspace{4cm}}}
\caption{Shared-curve residual attribution (RQ2) versus plot-specific trajectory attribution
(RQ3).}
\label{fig:attribution-comparison}
\end{figure}

\section{Reproducibility and appendix references}
\label{sec:reproducibility}

All split assignments, network weight initialisation, and Elastic Net/XGBoost cross-validation
folds use a fixed seed (42) unless explicitly swept as part of a stated robustness check (e.g. the
RQ1 winner-model reseed check, four additional seeds, reported in the Results chapter). Every
model reads its feature list from a single generated manifest
(\texttt{documentation/env\_feature\_sets\_manifest.csv}), so a reported set's exact column
membership can always be checked against that file rather than a hand-copied list.

The following are deliberately kept out of this chapter and referred to the Appendix:

\begin{itemize}
\setlength\itemsep{0pt}
\item \textbf{Appendix~\ref{app:hyperparameters}} -- full hyperparameter grids and selected
  values for every model; DNN/PINN architecture and training-loop settings (batch size, learning
  rate, optimiser, gradient clipping, early-stopping patience); the physics-weight ablation's
  exact settings.
\item \textbf{Appendix~\ref{app:feature-sets}} -- full Set1--Set5 column lists for all three
  research questions; per-variable dedup/VIF diagnostics; the complete correlation/VIF drop logs.
\item \textbf{Appendix~\ref{app:sensitivity}} -- secondary sensitivity settings not reported in
  the main text (dropout-rate and architecture-width sweeps; the narrow-gap temporal split's own
  results).
\end{itemize}

% TODO: verify -- the manifest-driven feature-set lookup (torch_data.py's ENV_TERRAIN_FEATURE_SETS)
% currently only covers RQ1; confirm whether RQ2/RQ3's own model-fitting scripts read the manifest
% directly or still require a separate wiring step before this chapter's Set1-5 description is
% fully accurate for all three research questions.
